% !TEX TS-program = pdflatex
% !TEX encoding = UTF-8 Unicode

% This file is a template using the "beamer" package to create slides for a talk or presentation
% - Giving a talk on some subject.
% - The talk is between 15min and 45min long.
% - Style is ornate.

% MODIFIED by Jonathan Kew, 2008-07-06
% The header comments and encoding in this file were modified for inclusion with TeXworks.
% The content is otherwise unchanged from the original distributed with the beamer package.

\documentclass{beamer}


% Copyright 2004 by Till Tantau <tantau@users.sourceforge.net>.
%
% In principle, this file can be redistributed and/or modified under
% the terms of the GNU Public License, version 2.
%
% However, this file is supposed to be a template to be modified
% for your own needs. For this reason, if you use this file as a
% template and not specifically distribute it as part of a another
% package/program, I grant the extra permission to freely copy and
% modify this file as you see fit and even to delete this copyright
% notice. 

\addtobeamertemplate{navigation symbols}{}{%
    \usebeamerfont{footline}%
    \usebeamercolor[fg]{footline}%
    \hspace{1em}%
    \insertframenumber/\inserttotalframenumber
}

\mode<presentation>
{
  \usetheme{Warsaw}
  % or ...

%  \setbeamercovered{transparent}
  % or whatever (possibly just delete it)
}

\usepackage{hyperref}

\usepackage[english]{babel}
% or whatever

\usepackage[utf8]{inputenc}
% or whatever

\usepackage{times}
\usepackage[T1]{fontenc}
% Or whatever. Note that the encoding and the font should match. If T1
% does not look nice, try deleting the line with the fontenc.



\title{Studying the Putnam}

\author[W.R. Casper] % (optional, use only with lots of authors)
{W.R. Casper\\\vspace{0.2in}
 {\tiny Problem Solving Seminar}}
% - Use the \inst{?} command only if the authors have different
%   affiliation.

\institute[California State University Fullerton] % (optional, but mostly needed)
{
  Department of Mathematics\\
  California State University Fullerton}
% - Use the \inst command only if there are several affiliations.
% - Keep it simple, no one is interested in your street address.

\subject{Talks}
% This is only inserted into the PDF information catalog. Can be left
% out. 



% If you have a file called "university-logo-filename.xxx", where xxx
% is a graphic format that can be processed by latex or pdflatex,
% resp., then you can add a logo as follows:

% \pgfdeclareimage[height=0.5cm]{university-logo}{university-logo-filename}
% \logo{\pgfuseimage{university-logo}}



% Delete this, if you do not want the table of contents to pop up at
% the beginning of each subsection:
\AtBeginSubsection[]
{
  \begin{frame}<beamer>{Outline}
    \tableofcontents[currentsection,currentsubsection]
  \end{frame}
}


% If you wish to uncover everything in a step-wise fashion, uncomment
% the following command: 

%\beamerdefaultoverlayspecification{<+->}


\begin{document}

\begin{frame}
  \titlepage
\end{frame}

%\begin{frame}{Outline}
  %\tableofcontents
  % You might wish to add the option [pausesections]
%\end{frame}

% Since this a solution template for a generic talk, very little can
% be said about how it should be structured. However, the talk length
% of between 15min and 45min and the theme suggest that you stick to
% the following rules:  

% - Exactly two or three sections (other than the summary).
% - At *most* three subsections per section.
% - Talk about 30s to 2min per frame. So there should be between about
%   15 and 30 frames, all told.

\begin{frame}{Introduction}
What is the Putnam Exam?
\pause
\begin{itemize}
\pause
\item running since $1938$
\pause
\item any undergrad can take it! (max $4$ years)
\pause
\item first Saturday of December
\pause
\item morning session: 8AM-11AM (6 problems)
\pause
\item afternoon session: 1PM-4PM (6 more problems)
\end{itemize}
\end{frame}

\begin{frame}{Benefits}
Why take the Putnam?
\begin{itemize}
\pause
\item bragging rights!
\pause
\item free lunch!
\pause
\item because I told you to
\pause
\item departmental scholarships
\pause
\item recommendation letters
\pause
\item not just for supergeniuses or the criminally insane
\end{itemize}
\end{frame}

\begin{frame}{Putnam practice}
Getting Putnam practice
\pause
\begin{itemize}
\pause
\item MAA Mathematics Magazine
\pause
\item College Math Journal
\pause
\item American Math Monthly
\pause
\item old problems / solutions: also in the Monthly!
\pause
\item Kiran Kedlaya's website
\end{itemize}
\end{frame}

\begin{frame}{Example problem}
\textbf{Putnam 2016 A1}:\\
Find the smallest integer $j$ such that for every polynomial $p(x)$ with integer coefficients
$$p^{(j)}(k)\ \ \text{is divisible by $2016$ for all integers $k$}.$$
\end{frame}

\begin{frame}{Strategy}
Strategies:
\pause
\begin{itemize}
\pause
\item carefully read the problem
\pause
\item underline critical info
\pause
\item classify the problem: what topic of math?  what useful results?
\pause
\item try to explore \emph{specific cases}
\end{itemize}
\end{frame}

\begin{frame}{Critical info}
\textbf{Putnam 2016 A1}:\\
Find the smallest integer $j$ such that for every polynomial $p(x)$ with integer coefficients
$$p^{(j)}(k)\ \ \text{is divisible by $2016$ for all integers $k$}.$$
\end{frame}

\begin{frame}{Critical info}
\textbf{Putnam 2016 A1}:\\
Find the smallest integer $j$ such that for every polynomial $p(x)$ with integer coefficients
$$p^{(j)}(k)\ \ \text{is {\color{red} divisible by $2016$} for all integers $k$}.$$
\end{frame}

\begin{frame}{Critical info}
\textbf{Putnam 2016 A1}:\\
Find the smallest integer $j$ such that for {\color{red}every polynomial} $p(x)$ with integer coefficients
$$p^{(j)}(k)\ \ \text{is divisible by $2016$ for all integers $k$}.$$
\end{frame}

\begin{frame}{Critical info}
\textbf{Putnam 2016 A1}:\\
Find the smallest integer $j$ such that for every polynomial $p(x)$ with {\color{red} integer coefficients}
$$p^{(j)}(k)\ \ \text{is divisible by $2016$ for all integers $k$}.$$
\end{frame}

\begin{frame}{Critical info}
\textbf{Putnam 2016 A1}:\\
Find the {\color{red} smallest integer} $j$ such that for every polynomial $p(x)$ with integer coefficients
$$p^{(j)}(k)\ \ \text{is divisible by $2016$ for all integers $k$}.$$
\end{frame}

\begin{frame}{Classify the problem}
Topics: polynomials, divisibility
\pause
important ideas/theorems
\begin{itemize}
\pause
\item primes, modular arithmetic
\pause
\item derivatives of polynomials
\pause
\item \textbf{Fundamental Theorem of Arithmetic}:\\
every positive integer may be written as a product of primes
$$n=p_1p_2\dots p_r$$
moreover, given another prime factorization
$$n=q_1q_2\dots q_s$$
then $r=s$ and the $p_j$'s and $q_j$'s are the same, after reordering
\end{itemize}
\pause
Also, polynomials are linear combinations of \textbf{monomials}: $x^n$
\end{frame}

\begin{frame}{Classify the problem}
Fundamental Theorem of Arithmetic:
$$2016 = 2^53^27$$
\pause
If $n$ is divisible by $2016$ then
\begin{itemize}
\pause
\item $n$ is divisible by $2^5$
\pause
\item $n$ is divisible by $3^2$
\pause
\item $n$ is divisible by $2^7$
\end{itemize}
\pause
Also, polynomials are linear combinations of \textbf{monomials}: $x^n$
\end{frame}

\begin{frame}{Scratch work}
\begin{align*}
f(x) &= x\\
f'(k) &= 1\\
f''(k) &= 0\\
f'''(k) &= 0\\
f''''(k) &= 0\\
\end{align*}
\end{frame}

\begin{frame}{Scratch work}
\begin{align*}
f(x) &= x^2\\
f'(k) &= 2k\\
f''(k) &= 2\\
f'''(k) &= 0\\
f''''(k) &= 0\\
\end{align*}
\end{frame}

\begin{frame}{Scratch work}
\begin{align*}
f(x) &= x^3\\
f'(k) &= 3k^2\\
f''(k) &= 6k\\
f'''(k) &= 6\\
f''''(k) &= 0\\
\end{align*}
\end{frame}

\begin{frame}{Scratch work}
\begin{align*}
f(x) &= x^4\\
f'(k) &= 4k^3\\
f''(k) &= 12k^2\\
f'''(k) &= 24k\\
f''''(k) &= 24\\
\end{align*}
\end{frame}

\begin{frame}{Scratch work}
\begin{align*}
f(x) &= x^{29}\\
f'(k) &= 29k^{28}\\
f''(k) &= 29\cdot 28k^{27}\\
f'''(k) &= 29\cdot 28\cdot 27 k^{26}\\
f''''(k) &= 29\cdot 28\cdot 27\cdot 26 k^{25}
\end{align*}
How long until divisible by 2016?
\end{frame}

\begin{frame}{Scratch work}
We need five $2$'s, two $3$'s and a $7$:
$$f^{(5)}(k) = 29\cdot 28\cdot 27\cdot 26\cdot 25 k^{24}$$
\pause
\begin{itemize}
\pause
\item $28$ has a $7$ and two $2$'s
\pause
\item $27$ has three $3$'s
\pause
\item $26$ has a $2$
\pause
\item $24$ has three $2$'s
\end{itemize}
\pause
Go one more to get a $2$:
$$f^{(6)}(k) = 29\cdot 28\cdot 27\cdot 26\cdot 25\cdot 24 k^{23}$$
\end{frame}

\begin{frame}{Scratch work}
Almost-worst case scenario: we might to do seven derivatives to get a $7$:
\pause
$$f(x) = x^{27}$$
\pause
$$f^{(6)}(k) = 27\cdot 26\cdot 25\cdot 24\cdot 23\cdot 22 k^{21}$$
\pause
$$f^{(7)}(k) = 27\cdot 26\cdot 25\cdot 24\cdot 23\cdot 22\cdot 21 k^{20}$$
\end{frame}

\begin{frame}{Scratch work}
Worst case scenario: we might have to do an eighth derivative to get enough $2$'s:
\pause
$$f(x) = x^{8\cdot 7-1} = x^{55}$$
\pause
$$f^{(7)}(k) = 55\cdot 54\cdot 53\cdot 52\cdot 51\cdot 50\cdot 49 k^{48}$$
\pause
\hfil This doesn't have enough $2$'s!
\pause
$$f^{(8)}(k) = 55\cdot 54\cdot 53\cdot 52\cdot 51\cdot 50\cdot 49\cdot 48  k^{47}$$
\hfil We guess that the \textbf{answer is 8}
\end{frame}

\begin{frame}{Writing up the solution}
\begin{itemize}
\pause
\item Good writing is important!
\pause
\item ALWAYS write the answer first
\pause
\item Carefully explain any unusual notation
\pause
\item Emphasize any theorems you are using
\pause
\item Emphasize key points
\pause
\item Break up your solution into manageable pieces with Lemmas
\pause
\item Double-check you answered the original question!
\end{itemize}
\end{frame}

\begin{frame}{Writing up the solution}
\textbf{Solution}:\\
\pause
The answer is $8$.\\
\pause
To see this, our primary tool is the Fundamental Theorem of Arithmetic, which says each integer has a unique prime factorization (up to reordering).\\
\pause
Since $2016 = 2^5\cdot 3^2\cdot 7$, this means divisibility by $2016$ is the same as divisibility by $2^5$, $3^2$, and $7$ all at once.\\
\pause
Now, for $f(x) = x^{55}$ we can calculate
$$f^{(7)}(1) = 55\cdot 54\cdot 53\cdot 52\cdot 51\cdot 50\cdot 49,$$
which is not divisible by $2016$, since it isn't divisible by $2^5$.\\
\pause
Therefore the answer is at least $8$.\\
\pause
To prove that $8$ is enough, we'll prove two Lemmas.\\
\end{frame}

\begin{frame}{Writing up the solution}
\textbf{Lemma}: If $n$ and $k$ are integers with $k>0$, then one of the integers\\
\hfil $\{n-0,n-1,n-2,\dots,n-k+1\}$\\
must be a multiple of $k$.
\pause
\textbf{Proof}:\\
The value of $n$ modulo $k$, ie. the remainder of $n/k$, is an integer $r$ with $0\leq r \leq n-1$.\\
Therefore $n-r$ will be $0$ modulo $k$, ie. it will be a multiple of $k$.
\end{frame}

\begin{frame}{Writing up the solution}
\textbf{Lemma}: If $n\geq 0$ is an integer and $f(x) = x^n$, then $f^{(j)}(k)$ is divisible by $2016$, for any integers $j$ and $k$ with $j\geq 8$.\\
\pause
\textbf{Proof}:\\
\pause
\hfil $f^{(8)}(x) = n(n-1)(n-2)(n-3)(n-4)(n-5)(n-6)(n-7)x^{n-8},$
\pause
and therefore for any integers $j$ and $k$ with $j\geq 8$, $f^{(j)}(k)$ will be a multiple of\\
\hfil $N = n(n-1)(n-2)(n-3)(n-4)(n-5)(n-6)(n-7).$
\pause
Consider the numbers
\hfil $\{n,n-1,n-2,n-3,n-4,n-5,n-6,n-7\}.$
\pause
At least four of the above numbers must be even, and therefore divisible by $2$. 
\pause
Actually, one must be a multiple of $4$!
\pause
Likewise, at least $2$ will be multiples of $3$. 
\pause
Finally, one of the numbers will be a multiple of $7$. 
\pause
Hence $N$ and $f^{(j)}(k)$ is divisible by $2016$.
\end{frame}

\begin{frame}{Writing up the solution}
To finish our solution, note that any polynomial will be of the form
$$f(x) = a_0 + a_1x + a_2x^2 + a_nx^n$$
for some integers $a_0,\dots, a_n$.\\
\pause
By our Lemma, when we take the $j$'th derivative for $j\geq 8$, each of the powers of $x$ becomes a multiple of $2016$.\\
\pause
This makes $f^{(j)}(k)$ a linear combination of multiples of $2016$, and therefore also a multiple of $2016$.
\end{frame}

\begin{frame}{Example problem}
\textbf{Try it yourself!}:\\
Find the smallest integer $j$ such that for every polynomial $p(x)$ with integer coefficients
$$p^{(j)}(k)\ \ \text{is divisible by $2024$ for all integers $k$}.$$
\end{frame}

\end{document}


